% ============================================================
\chapter*{Preface}
\addcontentsline{toc}{chapter}{Preface}
% ============================================================

\section*{Aims of the Course}

This \textbf{Probability Theory} course is aimed at second-year undergraduate students
in mathematics, computer science, or the physical sciences. It provides a rigorous
introduction to the mathematical theory of probability, grounded in the Kolmogorov
axioms, while maintaining strong probabilistic intuition throughout.

\medskip

Probability occupies a central place in modern mathematics and its applications.
From statistical physics to machine learning, from quantitative finance to molecular
biology, probabilistic modelling has become an indispensable tool for understanding
and quantifying uncertainty.

\section*{Prerequisites}

To approach this course effectively, students should be comfortable with:
\begin{itemize}
    \item \textbf{Real analysis:} sequences, series, integrals (Riemann), uniform
          convergence, basic dominated convergence results.
    \item \textbf{Linear algebra:} vector spaces, matrices, eigenvalues
          (needed for the Markov chains chapter).
    \item \textbf{Set theory:} set operations, countability, cardinality.
    \item \textbf{Combinatorics:} permutations, combinations, multiplication principle.
\end{itemize}

\section*{Course Organisation}

The course comprises twelve chapters, arranged in a logical progression:

\begin{description}
    \item[Chapters 1--2:] \textbf{Foundations.} We define the probability space
    $(\Omega, \mathcal{F}, \PP)$, state the Kolmogorov axioms, and develop
    conditional probability and independence. These two chapters establish the
    axiomatic framework on which everything else rests.

    \item[Chapters 3--5:] \textbf{Random variables.} We introduce discrete and then
    continuous random variables, with their distributions, cumulative distribution
    functions, and densities. Chapter~5 develops expectation, variance, and moments
    --- the fundamental tools of probabilistic calculation.

    \item[Chapters 6--8:] \textbf{Distributions and transformations.} Chapter~6
    presents a catalogue of standard distributions (binomial, Poisson, normal,
    exponential, gamma, etc.). Chapter~7 studies functions of random variables,
    and Chapter~8 introduces joint distributions, covariance, and correlation.

    \item[Chapters 9--11:] \textbf{Limit theorems.} This is the heart of the theory:
    the law of large numbers (weak and strong) in Chapter~9, the central limit
    theorem in Chapter~10, and characteristic functions in Chapter~11 --- a powerful
    tool that yields an elegant proof of the CLT.

    \item[Chapter 12:] \textbf{Introduction to Markov chains.} The final chapter
    opens the door to stochastic processes by studying Markov chains on finite or
    countable state spaces.
\end{description}

\section*{Methodology and Notation}

Throughout the course we adopt the following conventions:

\begin{itemize}
    \item \textbf{Definitions} are boxed and numbered. Every new concept is first
    defined rigorously before being illustrated by examples.

    \item \textbf{Theorems} are stated precisely and, wherever possible, proved in
    full. When a proof goes beyond the scope of the course, we provide a sketch or
    a reference.

    \item \textbf{Examples} are plentiful and detailed. They are an essential
    complement to the abstract theory.

    \item \textbf{Exercises} at the end of sections allow students to check their
    understanding and develop computational techniques.
\end{itemize}

\paragraph{Main notation.}
\begin{center}
\renewcommand{\arraystretch}{1.3}
\begin{tabular}{cl}
    \toprule
    \textbf{Symbol} & \textbf{Meaning} \\
    \midrule
    $\Omega$ & Sample space (set of all outcomes) \\
    $\mathcal{F}$ & $\sigma$-algebra (collection of events) \\
    $\PP$ & Probability measure \\
    $\E[X]$ & Expectation of the random variable $X$ \\
    $\Var(X)$ & Variance of $X$ \\
    $\Cov(X,Y)$ & Covariance of $X$ and $Y$ \\
    $F_X$ & Cumulative distribution function of $X$ \\
    $f_X$ & Probability density function of $X$ \\
    $\varphi_X(t)$ & Characteristic function of $X$ \\
    $X \sim \mathcal{L}$ & $X$ has distribution $\mathcal{L}$ \\
    $\xrightarrow{\text{a.s.}}$ & Almost sure convergence \\
    $\xrightarrow{\PP}$ & Convergence in probability \\
    $\xrightarrow{\mathcal{L}}$ & Convergence in distribution \\
    \bottomrule
\end{tabular}
\end{center}

\section*{Historical Perspective}

Probability theory has a rich history that deserves brief mention. The earliest
formal work dates to the correspondence between \textbf{Blaise Pascal} and
\textbf{Pierre de Fermat} in 1654 concerning the problem of points. In the
eighteenth century, \textbf{Jakob Bernoulli} proved the first version of the law
of large numbers in his \textit{Ars Conjectandi} (1713), while \textbf{Abraham
de Moivre} obtained a primitive form of the central limit theorem.

\medskip

In the nineteenth century, \textbf{Pierre-Simon de Laplace} synthesised the
probabilistic knowledge of his era in his \textit{Th\'eorie analytique des
probabilit\'es} (1812). \textbf{Pafnuty Chebyshev}, \textbf{Andrey Markov}, and
\textbf{Aleksandr Lyapunov} developed the analytic tools that would lead to the
modern limit theorems.

\medskip

The decisive revolution came in 1933, when \textbf{Andrey Kolmogorov} published his
\textit{Grundbegriffe der Wahrscheinlichkeitsrechnung}, founding probability theory
on measure theory. This axiomatisation, which we adopt throughout this course,
resolved numerous paradoxes and opened the way to the modern theory of stochastic
processes.

\section*{Advice to Students}

\begin{enumerate}
    \item \textbf{Work regularly.} Probability cannot be learnt the night before the
    exam. Each chapter builds on the previous ones.

    \item \textbf{Do the exercises.} Passive understanding of the theory is not
    enough. It is by solving problems that one develops probabilistic intuition.

    \item \textbf{Draw pictures.} Sketch sets, plot densities, draw distribution
    functions. Visualisation is a precious ally.

    \item \textbf{Check special cases.} When you obtain a general formula, test it
    on simple examples. This catches errors and strengthens understanding.

    \item \textbf{Distinguish discrete from continuous.} Many errors arise from
    confusing sums with integrals, probability mass functions with densities.
    Always be aware of the setting in which you are working.
\end{enumerate}

\section*{Main References}

This course draws on many classical texts, including:
\begin{itemize}
    \item J.~Jacod and P.~Protter, \textit{Probability Essentials}, Springer.
    \item P.~Billingsley, \textit{Probability and Measure}, Wiley.
    \item W.~Feller, \textit{An Introduction to Probability Theory and Its Applications},
          Vols.~I and~II, Wiley.
    \item G.~Grimmett and D.~Stirzaker, \textit{Probability and Random Processes}, Oxford.
    \item R.~Durrett, \textit{Probability: Theory and Examples}, Cambridge.
    \item J.~Rosenthal, \textit{A First Look at Rigorous Probability Theory}, World Scientific.
    \item S.~Ross, \textit{A First Course in Probability}, Pearson.
\end{itemize}

\bigskip

\hfill\textit{Happy reading, and good luck with your studies.}
